\chapter{备选：恩捷股份（002812）}

\begin{dashboardbox}[结论卡]
\textbf{现价：}47.84元 \quad \textbf{总市值：}469.46亿元 \quad \textbf{正式目标：}不发布 \quad \textbf{外部零权敏感性：}103元 \\
\textbf{基准/最终空间：}不发布 \quad \textbf{外部目标相对现价：}+115.3\% \quad \textbf{AStock翻倍概率：}不估计 \\
\textbf{核心判断：}盈利修复真实、现金流为正，且原始103元目标可审计；但目标建立在2027E利润跃升和单一券商的量价/利用率假设上，未通过正式估值资格门槛。
\end{dashboardbox}

\section{已经确认的修复}

恩捷股份2025年收入136.33亿元、归母净利润1.43亿元；2026Q1收入39.08亿元、归母2.60亿元、扣非2.61亿元，毛利率由2025年的18.77\%改善至28.09\%。H1预告归母7.36--9.00亿元、扣非7.62--9.30亿元，说明扭亏与毛利修复已进入报表。

现金质量相对锂盐和存储标的更好。FY2025经营现金流11.44亿元，2026Q1为2.07亿元，均为正。不过，隔膜业务资本开支高，新产线投放与利用率爬坡仍可能重新占用现金，不能仅以当前正现金流推断2027年盈利。

公司官方预告确认下游需求旺盛、产销量增长、价格较上年低点企稳回升、综合毛利率显著改善。官方没有披露H1出货量、客户涨价、实际单平利润、有效产能、利用率与良率，这些正是103元目标的关键变量。

\section{券商模型的强项与集中风险}

东吴证券预计2026年出货超过160亿平方米，2027年有效产能约220亿平方米，并假设湿法隔膜利用率从2026年的90\%提升至2027年的97\%；其2026/2027/2028归母净利润预测为22.79/50.68/62.01亿元，对应EPS 2.32/5.16/6.31元。目标价103元来自2027E EPS 5.16元乘20倍PE。

该模型的优点是经营桥清楚：出货平方米、单平价格、单位利润、利用率与折旧摊薄能够连接到EPS。缺点是关键参数高度同源。客户名称、涨价比例、订单、有效产能与新线良率没有公司级量化证据。一个完整的券商模型并不等于多个独立证据。

\begin{exhibitbox}[Exhibit 6：隔膜恢复桥的证据层级]
\small
\begin{tabularx}{\exhibitboxwidth}{L{3cm} L{3.4cm} L{3.4cm} X}
\toprule
变量 & 公司已披露 & 券商预测 & 估值处理 \\
\midrule
需求/价格 & 需求旺盛；价格企稳回升 & 部分大客户、海外与3C涨价仍在落地 & 方向给信用，幅度折价 \\
出货量 & 产销量增长，无吨位 & H1 76--77亿平；全年160亿平以上 & 不作为已实现事实 \\
单平利润 & 未披露 & Q2约0.15元；2027约0.20--0.25元 & Bull核心参数，需Q3验证 \\
利用率/产能 & 未披露同口径有效产能 & 2026/2027利用率90\%/97\%；2027有效产能220亿平 & 高度敏感，未验证前限制倍数 \\
现金流 & FY2025与Q1均为正 & 扩产资本开支仍较高 & 需要OCF覆盖扩产 \\
\bottomrule
\end{tabularx}
\sourcenote{公司公告；东吴证券2026年7月10日原始报告。}
\end{exhibitbox}

\section{估值：2027盈利必须被提前资本化}

当前价格对应2027E EPS约9.27倍。翻倍至95.68元需要18.54倍，接近券商20倍假设。也就是说，市场不仅要相信利润从2026年的22.79亿元增至2027年的50.68亿元，还要在2026年末前给予接近完整恢复周期的倍数。

东吴证券2027E EPS 5.16元乘6倍、13倍和20倍，可机械得到30.96元、67.08元和103.20元。三档全部依赖同一券商盈利分母，且公司尚未披露足以验证出货、单平利润、利用率、良率与客户涨价的参数，因此本报告不为其赋概率、不混合现价，也不发布正式目标。该区间只回答“若卖方分母兑现，估值可能落在哪里”，不构成AStock预期收益。

\section{为什么排名落后于中国船舶}

恩捷拥有显式目标，而中国船舶没有；但研究排名不等于目标价排名。中国船舶的订单与交付能见度、Q1现金流和利润质量更强，经营桥中的公司级缺口更少。恩捷的客户涨价、利用率、良率与单平利润尚未公开量化，因此在核心资格上必须让位于证据质量。

\begin{riskbox}[降级条件]
若Q3单平利润低于约0.12元、客户涨价继续延迟、新线良率或利用率不及预期、经营现金流被应收与资本开支吞噬，103元牛市情景应删除。只有出货约40亿平以上、单平利润约0.15元、客户涨价与现金流共同验证，才可重审核心资格。
\end{riskbox}

\clearpage
